\chapter{Eccezioni, Asserzioni e Logging}

\section{Gestione delle Eccezioni}

Il passaggio dalla programmazione accademica a quella professionale coincide con la consapevolezza che il codice deve operare in un ambiente ostile (dati errati, bug, risorse esterne indisponibili).

In presenza di un'anomalia, il software deve garantire tre requisiti minimi:
\begin{enumerate}
    \item \textbf{Notifica}: Informare l'utilizzatore del problema occorso.
    \item \textbf{Persistenza}: Salvare lo stato corrente per evitare la perdita di dati.
    \item \textbf{Terminazione Controllata}: Consentire una chiusura del programma che rilasci correttamente le risorse impegnate.
\end{enumerate}
\vspace{\baselineskip}

Storicamente, i metodi segnalavano il fallimento tramite valori sentinella (\texttt{null}, \texttt{-1}, \texttt{false}). Questo approccio è oggi considerato \textbf{pericoloso} e \textbf{poco informativo}.

Confrontiamo l'evoluzione delle API Java per capire il vantaggio delle eccezioni:

\begin{itemize}
    \item \textbf{Vecchio Stile (Legacy)}: \texttt{file.delete()} restituisce \texttt{false} se fallisce. 
    \begin{itemize}
        \item \textit{Problema}: Perché e' fallito? Permessi negati? File inesistente? Disco protetto? Non lo sapremo mai.
    \end{itemize}
    \item \textbf{Nuovo Stile (NIO)}: \texttt{Files.delete(path)} lancia un'eccezione, ad esempio: 
    \begin{center}
        \texttt{NoSuchFileException}, \texttt{AccessDeniedException}, ecc...
    \end{center}
    
    \begin{itemize}
        \item \textit{Vantaggio}: L'errore e' auto-esplicativo e costringe il programmatore a decidere come gestirlo.
    \end{itemize}
\end{itemize}
\vspace{\baselineskip}

Come vedremo presto, qui entra in gioco il meccanismo di \textit{throw}; quando un metodo lancia un'eccezione:
\begin{enumerate}
    \item Non viene restituito alcun valore al chiamante.
    \item Il controllo passa immediatamente al meccanismo di gestione delle eccezioni della JVM.
    \item La JVM cerca un \textit{handler} compatibile risalendo la catena delle chiamate (\textit{unwinding the stack}).
\end{enumerate}

\begin{tcolorbox}[colback=blue!5!white,colframe=blue!75!black,title=Best Practice: Eccezioni vs Performance]
\textbf{1. Usa le eccezioni solo dove necessario}: Non utilizzare mai le eccezioni per gestire il normale flusso di controllo del programma (es. per uscire da un loop). Creare un'eccezione è "costoso" per la JVM perché deve catturare l'intero \textit{stack trace} (lo stato di tutte le chiamate ai metodi in quel momento).

\textbf{2. Preferisci metodi che lanciano eccezioni}: Quando esplori le API di Java, preferisci sempre i metodi che segnalano errori tramite eccezioni rispetto a quelli che restituiscono codici di errore (come \texttt{null} o \texttt{-1}).
\begin{itemize}
\item Un'eccezione \textbf{non può essere ignorata}: se non la gestisci, il programma si ferma, impedendo al bug di propagarsi silenziosamente.
\item Un codice di errore richiede un \texttt{if} manuale: se il programmatore si dimentica di scriverlo, il programma continuerà a girare in uno stato instabile.
\end{itemize}
\end{tcolorbox}

\subsection{La Gerarchia delle Eccezioni}

In Java, la gestione degli errori è organizzata in una gerarchia rigorosa. Di seguito è riportata la struttura standard UML, ottimizzata per la leggibilità e la precisione dei collegamenti.

\begin{figure}[ht]
\centering
\begin{tikzpicture}[
    node distance=0.4cm and 1.5cm,
    >= {Triangle[open, length=7pt, width=7pt]}, % Freccia UML standard (vuota)
    % Stili dei blocchi
    base/.style={rectangle, rounded corners=1pt, draw=none, thick, minimum width=3.8cm, minimum height=0.7cm, align=center, font=\sffamily\bfseries\scriptsize, drop shadow},
    pkg/.style={base, fill=magenta, text=white, minimum width=2.5cm},
    root/.style={base, fill=blue!30!black, text=white},
    checked/.style={base, fill=green!45!black, text=white},
    runtime_main/.style={base, fill=blue!60!black, text=white},
    unchecked/.style={base, fill=cyan!70!black, text=white},
    error_node/.style={base, fill=red!70!black, text=white}
]

% --- LIVELLO 1: Radice ---
\node[pkg] (java) {java.lang};
\node[root, below=0.3cm of java] (object) {Object};
\node[root, below=0.6cm of object] (throwable) {Throwable};

% --- LIVELLO 2: Rami Principali ---
\node[root, below left=1.2cm and 1.0cm of throwable] (exception) {Exception};
\node[root, below right=1.2cm and 1.0cm of throwable] (error) {Error};

% --- LIVELLO 3: Figli di Exception (Checked) ---
\node[checked, below=0.8cm of exception, xshift=0.0cm] (ioe) {IOException};
\node[checked, below=0.2cm of ioe] (cnf) {ClassNotFoundException};
\node[checked, below=0.2cm of cnf] (sql) {SQLException};
\node[runtime_main, below=0.2cm of sql] (runtime) {RuntimeException};

% --- LIVELLO 4: Figli di RuntimeException ---
\node[unchecked, below=0.8cm of runtime, xshift=0.0cm] (arith) {ArithmeticException};
\node[unchecked, below=0.2cm of arith] (npe) {NullPointerException};
\node[unchecked, below=0.2cm of npe] (nfe) {NumberFormatException};
\node[unchecked, below=0.2cm of nfe] (iobe) {IndexOutOfBoundsException};

% --- LIVELLO 5: Figli di IndexOutOfBounds ---
\node[unchecked, below=0.8cm of iobe, xshift=0.0cm] (aiobe) {ArrayIndexOutOfBounds...};
\node[unchecked, below=0.2cm of aiobe] (siobe) {StringIndexOutOfBounds...};

% --- LIVELLO 3 (Bis): Figli di Error ---
\node[error_node, below=0.8cm of error, xshift=0.0cm] (vme) {VirtualMachineError};
\node[error_node, below=0.2cm of vme] (oom) {OutOfMemoryError};
\node[error_node, below=0.2cm of oom] (soe) {StackOverflowError};

% --- CONNESSIONI (Logica UML Rigorosa) ---

% 1. Throwable -> Object
\draw[->, thick] (throwable.north) -- (object.south);

% 2. Exception & Error -> Throwable (Connessione a T)
\coordinate (T_joint) at ([yshift=-0.5cm]throwable.south);
\draw[thick] (exception.north) |- (T_joint);
\draw[thick] (error.north) |- (T_joint);
\draw[->, thick] (T_joint) -- (throwable.south);

% 3. Famiglia Exception (Linee a cascata)
\coordinate (E_spine) at ([xshift=-0.5cm]ioe.west);
\draw[thick] (E_spine |- ioe.west) -- (E_spine |- runtime.west); % Spina verticale
\draw[->, thick] (E_spine |- ioe.west) |- (exception.west);     % Attacco al genitore
\draw[thick] (ioe.west) -- (E_spine |- ioe.west);
\draw[thick] (cnf.west) -- (E_spine |- cnf.west);
\draw[thick] (sql.west) -- (E_spine |- sql.west);
\draw[thick] (runtime.west) -- (E_spine |- runtime.west);

% 4. Famiglia RuntimeException
\coordinate (R_spine) at ([xshift=-0.5cm]arith.west);
\draw[thick] (R_spine |- arith.west) -- (R_spine |- iobe.west);
\draw[->, thick] (R_spine |- arith.west) |- (runtime.west);
\draw[thick] (arith.west) -- (R_spine |- arith.west);
\draw[thick] (npe.west) -- (R_spine |- npe.west);
\draw[thick] (nfe.west) -- (R_spine |- nfe.west);
\draw[thick] (iobe.west) -- (R_spine |- iobe.west);

% 5. Famiglia IndexOutOfBounds
\coordinate (I_spine) at ([xshift=-0.5cm]aiobe.west);
\draw[thick] (I_spine |- aiobe.west) -- (I_spine |- siobe.west);
\draw[->, thick] (I_spine |- aiobe.west) |- (iobe.west);
\draw[thick] (aiobe.west) -- (I_spine |- aiobe.west);
\draw[thick] (siobe.west) -- (I_spine |- siobe.west);

% 6. Famiglia Error
\coordinate (Err_spine) at ([xshift=-0.5cm]vme.west);
\draw[thick] (Err_spine |- vme.west) -- (Err_spine |- soe.west);
\draw[->, thick] (Err_spine |- vme.west) |- (error.west);
\draw[thick] (vme.west) -- (Err_spine |- vme.west);
\draw[thick] (oom.west) -- (Err_spine |- oom.west);
\draw[thick] (soe.west) -- (Err_spine |- soe.west);

\end{tikzpicture}
\caption{Gerarchia delle eccezioni in Java (Dettaglio UML)}
\label{fig:eccezioni_vettoriale_uml}
\end{figure}


